\begin{thebibliography}{39}
\bibitem[1]{ref1} M. Ajtai. The complexity of the pigeonhole principle. Combinatorica, 14(4):417–433, 1994.
\bibitem[2]{ref2} M. Alekhnovich, E. Ben-Sasson, A. A. Razborov, and A. Wigderson. Pseudorandom generators in propositional proof complexity. SIAM Journal on Computing, 34(1):67–88, 2004.
\bibitem[3]{ref3} E. Ben-Sasson and R. Impagliazzo. Random CNF’s are hard for the polynomial calculus. Computational Complexity, 19(4):501–519, 2010.
\bibitem[4]{ref4} M. L. Bonet, C. Domingo, R. Gavaldà, A. Maciel, and T. Pitassi. Non-automatizability of bounded-depth Frege proofs. Computational Complexity, 13(1-2):47–68, 2004.
\bibitem[5]{ref5} S. R. Buss. Bounded arithmetic. Studies in Proof Theory. Lecture Notes, 3. Napoli: Bibliopolis. VII, 221 p. (1986)., 1986.
\bibitem[6]{ref6} R. Fagin. Generalized first-order spectra and polynomial-time recognizable sets. Complexity of Comput., Proc. Symp. appl. Math., New York City 1973, 43-73 (1974)., 1974.
\bibitem[7]{ref7} E. Jeřábek. Dual weak pigeonhole principle, Boolean complexity, and derandomization. Annals of Pure and Applied Logic, 129(1-3):1–37, 2004.
\bibitem[8]{ref8} E. Jeřábek. Approximate counting in bounded arithmetic. The Journal of Symbolic Logic, 72(3):959–993, 2007.
\bibitem[9]{ref9} J. Krajíček. Bounded arithmetic, propositional logic, and complexity theory, volume 60. Cambridge: Cambridge Univ. Press, 1995.
\bibitem[10]{ref10} J. Krajíček. Interpolation theorems, lower bounds for proof systems, and independence results for bounded arithmetic. The Journal of Symbolic Logic, 62(2):457–486, 1997.
\bibitem[11]{ref11} J. Krajíček. On the weak pigeonhole principle. Fundamenta Mathematicae, 170(1-2):123–140, 2001.
\bibitem[12]{ref12} J. Krajíček. Tautologies from pseudo-random generators. The Bulletin of Symbolic Logic, 7(2):197–212, 2001.
\bibitem[13]{ref13} J. Krajíček. Diagonalization in proof complexity. Fundamenta Mathematicae, 182(2):181–192, 2004.
\bibitem[14]{ref14} J. Krajíček. Dual weak pigeonhole principle, pseudo-surjective functions, and provability of circuit lower bounds. The Journal of Symbolic Logic, 69(1):265–286, 2004.
\bibitem[15]{ref15} J. Krajíček. Structured pigeonhole principle, search problems and hard tautologies. The Journal of Symbolic Logic, 70(2):616–630, 2005.
\bibitem[16]{ref16} J. Krajíček. A proof complexity generator. In Logic, methodology and philosophy of science. Proceedings of the 13th international congress, Beijing, China, August 2007, pages 185–190. London: College Publications, 2009.
\bibitem[17]{ref17} J. Krajíček. A form of feasible interpolation for constant depth Frege systems. The Journal of Symbolic Logic, 75(2):774–784, 2010.
\bibitem[18]{ref18} J. Krajíček. Forcing with random variables and proof complexity, volume 382. Cambridge: Cambridge University Press, 2011.
\bibitem[19]{ref19} J. Krajíček. On the proof complexity of the Nisan-Wigderson generator based on a hard NP $\cap$ coNP function. Journal of Mathematical Logic, 11(1):11–27, 2011.
\bibitem[20]{ref20} J. Krajíček. A reduction of proof complexity to computational complexity for AC0 [p] Frege systems. Proceedings of the American Mathematical Society, 143(11):4951–4965, 2015.
\bibitem[21]{ref21} J. Krajíček. Proof complexity, volume 170. Cambridge: Cambridge University Press, 2019.
\bibitem[22]{ref22} J. Krajíček. Small Circuits and Dual Weak PHP in the Universal Theory of P-Time Algorithms. ACM Transactions on Computational Logic, 22(2), 2021.
\bibitem[23]{ref23} J. Krajíček and P. Pudlák. Propositional proof systems, the consistency of first order theories and the complexity of computations. The Journal of Symbolic Logic, 54(3):1063–1079, 1989.
\bibitem[24]{ref24} J. Krajíček and P. Pudlák. Some consequences of cryptographical conjectures for S21 and EF. Information and Computation, 140(1):82–94, 1998.
\bibitem[25]{ref25} J. Krajíček, P. Pudlák, and A. Woods. An exponential lower bound to the size of bounded depth Frege proofs of the pigeonhole principle. Random Structures \& Algorithms, 7(1):15–39, 1995.
\bibitem[26]{ref26} N. Nisan and A. Wigderson. Hardness vs randomness. Journal of Computer and System Sciences, 49(2):149–167, 1994.
\bibitem[27]{ref27} J. Paris and A. Wilkie. Counting problems in bounded arithmetic. Methods in mathematical logic, Proc. 6th Latin Amer. Symp., Caracas/Venez. 1983, Lect. Notes Math. 1130, 317-340 (1985)., 1985.
\bibitem[28]{ref28} J. Pich. Nisan-Wigderson generators in proof systems with forms of interpolation. Mathematical Logic Quarterly (MLQ), 57(4):379–383, 2011.
\bibitem[29]{ref29} J. Pich. Circuit lower bounds in bounded arithmetics. Annals of Pure and Applied Logic, 166(1):29–45, 2015.
\bibitem[30]{ref30} J. Pich. Logical strength of complexity theory and a formalization of the PCP theorem in bounded arithmetic. Logical Methods in Computer Science, 11(2):38, 2015.
\bibitem[31]{ref31} J. Pich. Learning algorithms from circuit lower bounds. arXiv, 2020. arXiv:2012.14095.
\bibitem[32]{ref32} J. Pich and R. Santhanam. Strong Co-Nondeterministic Lower Bounds for NP Cannot Be Proved Feasibly. In Proceedings of the 53rd Annual ACM SIGACT Symposium on Theory of Computing, STOC 2021, pages 223–233, 2021.
\bibitem[33]{ref33} T. Pitassi, P. Beame, and R. Impagliazzo. Exponential lower bounds for the pigeonhole principle. Computational Complexity, 3(2):97–140, 1993.
\bibitem[34]{ref34} P. Pudlák. Lower bounds for resolution and cutting plane proofs and monotone computations. The Journal of Symbolic Logic, 62(3):981–998, 1997.
\bibitem[35]{ref35} P. Pudlák. Incompleteness in the finite domain. The Bulletin of Symbolic Logic, 23(4):405–441, 2017.
\bibitem[36]{ref36} P. Pudlák. The canonical pairs of bounded depth Frege systems. Annals of Pure and Applied Logic, 172(2):42, 2021. Id/No 102892.
\bibitem[37]{ref37} A. A. Razborov. Pseudorandom generators hard for k-DNF resolution and polynomial calculus resolution. Annals of Mathematics. Second Series, 181(2):415–472, 2015.
\bibitem[38]{ref38} D. Sokolov. Pseudorandom Generators, Resolution and Heavy Width. In S. Lovett, editor, 37th Computational Complexity Conference (CCC 2022), volume 234 of Leibniz International Proceedings in Informatics (LIPIcs), pages 15:1–15:22, Dagstuhl, Germany, 2022. Schloss Dagstuhl – Leibniz-Zentrum für Informatik. doi:10.4230/LIPIcs.CCC.2022.15.
\bibitem[39]{ref39} A. R. Woods. Approximating the structures accepted by a constant depth circuit or satisfying a sentence – a nonstandard approach. In Logic and random structures. DIMACS workshop, November 5–7, 1995, pages 109–130. DIMACS/AMS, 1997.
\end{thebibliography}
